% =====================================================================
%  Thème 2 — Angles associés — NIVEAU MOYEN
%  Exercices
% =====================================================================
\documentclass[11pt,a4paper]{article}
\usepackage[utf8]{inputenc}
\usepackage[T1]{fontenc}
\usepackage{lmodern}
\usepackage[french]{babel}
\usepackage{amsmath,amssymb,mathtools}
\usepackage{icomma}
\usepackage{xcolor}
\usepackage{colortbl}
\usepackage{tikz}
\usepackage[most]{tcolorbox}
\usepackage{enumitem}
\usepackage{array}
\usepackage[a4paper,margin=2.2cm]{geometry}
\usetikzlibrary{arrows.meta,calc}

\definecolor{primary}{RGB}{214,51,108}
\definecolor{primarydark}{RGB}{140,20,64}

\DeclareMathOperator{\tg}{tg}
\DeclareMathOperator{\cotg}{cotg}
\newcommand{\degr}{^{\circ}}
\newcommand{\Z}{\mathbb{Z}}

\newtcolorbox{consigne}{enhanced,breakable,colback=primary!5,colframe=primary,
  boxrule=0.8pt,arc=2pt,left=3mm,right=3mm,top=2mm,bottom=2mm}

\newcounter{exo}
\newcommand{\exo}[1]{\medskip\refstepcounter{exo}%
  \noindent{\large\bfseries\color{primarydark}Exercice \theexo}%
  \ \textcolor{primary}{\textbullet}\ \textbf{#1}\par\nopagebreak\smallskip}

\setlength{\parindent}{0pt}
\pagestyle{empty}

\begin{document}

\begin{center}
{\LARGE\bfseries\color{primary}Thème 2 — Angles associés}\\[2pt]
{\color{primarydark}\textsc{Niveau Moyen — Exercices}}
\end{center}
\hrule height 1pt
\medskip

\begin{consigne}
\textbf{Objectif :} enchaîner \emph{deux} étapes — reconnaître la famille d'angles associés, puis
conclure (valeur exacte ou expression simplifiée). \textbf{Sans calculatrice.} Nomme la famille utilisée.
\end{consigne}

\exo{Réduire, puis évaluer}
Écris l'angle sous la forme $\pi\pm\beta$, $\frac{\pi}{2}\pm\beta$ ou $-\beta$ avec $\beta$ remarquable,
puis donne la valeur exacte.
\[
\textbf{(a)}\ \sin 210\degr \qquad
\textbf{(b)}\ \cos(-120\degr) \qquad
\textbf{(c)}\ \tg 225\degr \qquad
\textbf{(d)}\ \cos 135\degr \qquad
\textbf{(e)}\ \sin\frac{5\pi}{6}
\]

\exo{Simplifier une expression littérale}
Simplifie chaque expression le plus possible (le résultat est un nombre, $\pm\sin\alpha$, $\pm\cos\alpha$\dots).
\[
\textbf{(a)}\ \cos\!\left(\tfrac{\pi}{2}+\alpha\right)+\sin(\pi-\alpha) \qquad
\textbf{(b)}\ \sin(-\alpha)+\sin(\pi+\alpha)
\]
\[
\textbf{(c)}\ \dfrac{\tg(\pi+\alpha)}{\tg(-\alpha)} \qquad
\textbf{(d)}\ \cos(\pi-\alpha)+\cos(-\alpha)
\]

\exo{Lesquelles valent $\sin 5\degr$ ?}
On pose $\alpha=5\degr$ (angle du 1\textsuperscript{er} quadrant). Parmi les expressions suivantes,
détermine \emph{lesquelles} sont égales à $\sin 5\degr$, en justifiant par la famille d'angles.
\[
\textbf{(a)}\ \sin 175\degr \qquad
\textbf{(b)}\ -\cos 95\degr \qquad
\textbf{(c)}\ \sin(-5\degr) \qquad
\textbf{(d)}\ -\sin 185\degr \qquad
\textbf{(e)}\ \cos(-85\degr)
\]

\exo{Lesquelles valent $\tg 10\degr$ ?}
On pose $\alpha=10\degr$. Détermine lesquelles de ces expressions sont égales à $\tg 10\degr$.
\[
\textbf{(a)}\ \tg 190\degr \qquad
\textbf{(b)}\ \cotg 80\degr \qquad
\textbf{(c)}\ \tg 170\degr \qquad
\textbf{(d)}\ -\tg(-10\degr)
\]

\exo{Vrai ou Faux}
Justifie brièvement (nomme la bonne famille).
\begin{enumerate}[label=\textbf{(\alph*)},leftmargin=2em,itemsep=1pt]
\item $5\degr$ et $175\degr$ sont anti-supplémentaires.
\item $5\degr$ et $85\degr$ sont complémentaires.
\item $\cos 100\degr=-\sin 10\degr$.
\item $\sin 250\degr=\sin 70\degr$.
\end{enumerate}

\end{document}
